\section{Introduction}
The proliferation of Generative Artificial Intelligence (GenAI) has fundamentally altered the landscape of religious education and personal study. However, most existing AI tools function primarily as ``Search Engines with a Voice,'' providing deductive answers that often prematurely terminate the user's own reflective process. This phenomenon, which we term the ``Mirror Effect'' or ``Stochastic Parrot'' response \cite{bender2021parrots}, occurs when the AI reflects the user's cognitive biases or limits its feedback to a summary of the provided input, thereby offering zero ``Value-Add'' to the serious scholar.

\subsection{The Problem of Shallow Engagement}
The risks of AI in theological contexts extend beyond mere inaccuracy. As recent scholarship on AI in theological education warns, the ease of generating summaries and outlines may circumvent the ``formative struggle'' that is essential for intellectual and spiritual maturity. The process of wrestling with complex texts---what Ricoeur \cite{ricoeur1976} calls the ``surplus of meaning''---is itself the transformative act. When an AI short-circuits this struggle by providing a pre-digested answer, it robs the user of the very encounter that produces growth.

\subsection{Computational Hermeneutics}
To address these limitations, we draw upon the emerging field of Computational Hermeneutics \cite{kommers2026}. This discipline posits that AI should be used as an interpretive ``accelerator'' that navigates the situatedness and plural meanings of ancient texts. Rather than serving as an oracle, the AI acts as a mid-wife (\textit{maieutics}) to help the user give birth to new insights that were previously hidden by linguistic or cultural barriers. Gadamer's \cite{gadamer1960} concept of the ``fusion of horizons'' is instructive here: true understanding occurs not when the AI delivers a historical fact, but when the user's present-day ``horizon'' merges with the ancient text's ``horizon'' to produce a new, shared context of meaning.

\subsection{The Relational Disconnect}
A significant gap in contemporary Bible study literature is the ``Methodological Trap,'' where the study of sacred texts becomes a mechanical ritual of ``doing'' rather than ``relating'' \cite{packer1973}. Traditional Inductive Bible Study (IBS) methods, as formalized by Traina \cite{traina1952} and refined by Bauer and Traina \cite{bauer2011}, provide a rigorous three-step process of Observation, Interpretation, and Application (OIA). While invaluable, these methods can inadvertently lead to a ``checklist'' mentality that distances the reader from the Divine. In contrast, a relational hermeneutic rooted in the Father-Child paradigm suggests that the primary goal of the text is to facilitate a natural, spontaneous connection \cite{nih2024}. This paper proposes a framework that bridges this gap, using high-level computational tools to foster a child-like curiosity and intimacy.

\subsection{Paper Organization}
The remainder of this paper is organized as follows. Section 2 presents a gap analysis identifying the specific deficiencies in current AI-augmented Bible study tools. Section 3 establishes the theological foundation of the Relational Priority. Section 4 details the three-layer Scriptorium Cycle methodology. Section 5 specifies the implementation rules and discusses future research directions.
